\chapter{1977年河南卷}

\begin{problem}
化简下式：\( \left( 1-c^2 \right)^{- \frac{1}{2}}-\left\{ \left[ \left( 1+c \right)\left( 1-c \right) \right]^{\frac{1}{2}}+c^2\left( \left( 1+c \right)\left( 1-c \right) \right)^{- \frac{1}{2}} \right\} \).
\end{problem}

\begin{solution}
由原式中幂为负的项可知，必须有 \(1-c^2>0\)，即 \(-1<c<1\). 在此条件下，
\(\left[(1+c)(1-c)\right]^{\frac12}=(1-c^2)^{\frac12}\)，所以原式为
\(\left(1-c^2\right)^{-\frac12}-\left(1-c^2\right)^{\frac12}-c^2\left(1-c^2\right)^{-\frac12}\). 将同类项合并，得
\[
\left(1-c^2\right)^{-\frac12}\left[1-\left(1-c^2\right)-c^2\right]=0
\]
因此原式的值为 \(0\).
\end{solution}

\begin{problem}
求 \( \frac{\lg 3+\lg 2}{\frac{1}{4}\lg 16+\sqrt{\left( \lg 1-\lg 10 \right)^2}+\frac{1}{2}\lg 0.09} \) 的值.
\end{problem}

\begin{solution}
因为 \(\lg16=4\lg2\)，\(\lg1=0\)，\(\lg10=1\)，且
\(0.09=0.3^2\)，所以
\(\frac14\lg16=\lg2\)，\(\sqrt{(\lg1-\lg10)^2}=\sqrt{(-1)^2}=1\)，以及
\(\frac12\lg0.09=\lg0.3=\lg3-\lg10=\lg3-1\). 因此原式等于
\[
\frac{\lg3+\lg2}{\lg2+1+\lg3-1}
=\frac{\lg3+\lg2}{\lg3+\lg2}=1
\]
这里分母为 \(\lg6>0\)，约分合法.
\end{solution}

\begin{problem}
求等式 \( 5^{x^2-2x}=125 \) 中 \( x \) 的值.
\end{problem}

\begin{solution}
因为底数 \(5>0\) 且 \(5\ne1\)，指数函数 \(5^t\) 在实数范围内是单调函数，故由
\(5^{x^2-2x}=125=5^3\) 可得 \(x^2-2x=3\). 于是
\(x^2-2x-3=0\)，\(\left(x-3\right)\left(x+1\right)=0\).
所以 \(x=3\) 或 \(x=-1\). 代回原方程，分别有 \(5^{9-6}=5^3\) 和 \(5^{1+2}=5^3\)，两值均成立，故解为 \(x_1=3\)，\(x_2=-1\).
\end{solution}

\begin{problem}
已知 \( a \)，\( b \)，\( c \) 分别表示 \( \triangle ABC \) 的角 \(A\)，\(B\)，\(C\) 对边的长，求证：\( a\left( \sin B-\sin C \right)+b\left( \sin C-\sin A \right)+c\left( \sin A-\sin B \right)=0 \).
\end{problem}

\begin{solution}
由正弦定理，\(\frac{a}{\sin A}=\frac{b}{\sin B}=\frac{c}{\sin C}\). 交叉相乘可得
\(a\sin B=b\sin A\)，\(a\sin C=c\sin A\)，\(b\sin C=c\sin B\). 因此
\[
\begin{gathered}
a\left( \sin B-\sin C \right)+b\left( \sin C-\sin A \right)+c\left( \sin A-\sin B \right) \\
=a\sin B-a\sin C+b\sin C-b\sin A+c\sin A-c\sin B \\
=b\sin A-c\sin A+c\sin B-b\sin A+c\sin A-c\sin B=0
\end{gathered}
\]
上式正是所要证明的等式，故命题成立.
\end{solution}

\begin{problem}
如图，已知正方形 \(ABCD\) 的边 \(CD\) 上任意一点 \(E\)，延长 \(BC\) 到 \(F\)，使 \( CF=CE \)，设 \(BE\) 与 \(DF\) 相交于 \(G\)，求证 \( BG \perp DF \).

\bitmapfigure[max width=5.2cm]{img/1977/henan/q5_fig1.png}
\end{problem}

\begin{solution}
设正方形边长为 \(s>0\)，取直角坐标系使
\(B=(0,0)\)，\(C=(s,0)\)，\(D=(s,s)\). 设 \(CE=t\)，则 \(0\leq t\leq s\)，从而
\(E=(s,t)\). 由于 \(F\) 在 \(BC\) 的延长线上且 \(CF=CE=t\)，有 \(F=(s+t,0)\).

直线 \(BE\) 和 \(DF\) 的方向向量分别可取为
\(\overrightarrow{BE}=(s,t)\) 和 \(\overrightarrow{DF}=(t,-s)\). 它们的内积为
\[
\overrightarrow{BE}\mathbin{\cdot}\overrightarrow{DF}=st-ts=0
\]
因此 \(BE\perp DF\). 又因为 \(G\) 在 \(BE\) 和 \(DF\) 上，\(BG\) 与 \(BE\) 共线，所以 \(BG\perp DF\). 当 \(t=0\) 或 \(t=s\) 时，以上方向向量仍然成立，分别对应 \(G=C\) 或 \(G=D\)，结论同样成立.
\end{solution}

\begin{problem}
求证：\( \tan^2\alpha-\cot^2\alpha=-\frac{2\sin 4\alpha}{\sin^3 2\alpha} \).
\end{problem}

\begin{solution}
在等式有意义的条件下，\(\sin\alpha\ne0\) 且 \(\cos\alpha\ne0\)，因而 \(\sin2\alpha\ne0\). 由平方差公式和二倍角公式，左边为
\[
\tan^2\alpha-\cot^2\alpha
=\frac{\sin^2\alpha}{\cos^2\alpha}-\frac{\cos^2\alpha}{\sin^2\alpha}
=\frac{\sin^4\alpha-\cos^4\alpha}{\sin^2\alpha\cos^2\alpha}
=\frac{\sin^2\alpha-\cos^2\alpha}{\sin^2\alpha\cos^2\alpha}
\]
又因为 \(\sin^2\alpha-\cos^2\alpha=-\cos2\alpha\)，以及
\(\sin^2\alpha\cos^2\alpha=\frac14\sin^22\alpha\)，所以
再利用 \(\sin4\alpha=2\sin2\alpha\cos2\alpha\)，有
\[
\tan^2\alpha-\cot^2\alpha
=-\frac{4\cos2\alpha}{\sin^22\alpha}
=-\frac{4\cos2\alpha\sin2\alpha}{\sin^32\alpha}
=-\frac{2\sin4\alpha}{\sin^32\alpha}
\]
这正是右边，故等式成立.
\end{solution}

\begin{problem}
已知长方体的对角线长为 \(L\)，它与底面所成的角为 \( \alpha \)，底面两条对角线的夹角为 \( \beta \)，求长方体的体积.

\bitmapfigure[max width=5.8cm]{img/1977/henan/q7_fig1.png}
\end{problem}

\begin{solution}
设长方体的高为 \(AA_1\)，底面两条对角线为 \(AC\) 和 \(BD\). 空间对角线在底面上的射影是底面对角线 \(AC\)，所以在直角三角形 \(\triangle AA_1C\) 中，
\(AA_1=L\sin\alpha\)，\(AC=L\cos\alpha\).

矩形的两条对角线相等，且它们的夹角为 \(\beta\)，因此底面面积等于两条对角线所构成的四个三角形面积之和，即
\[
S=\frac12AC\cdot BD\sin\beta
=\frac12AC^2\sin\beta
=\frac12L^2\cos^2\alpha\sin\beta
\]
所以长方体的体积为
\[
V=AA_1\cdot S
=L\sin\alpha\cdot\frac12L^2\cos^2\alpha\sin\beta
=\frac12L^3\cos^2\alpha\sin\alpha\sin\beta
\]
\end{solution}

\begin{problem}
某工厂科研小组，对一项生产工艺过程总结出产量指标函数和消耗函数为：\( f_1(x)=a_1x^2+\frac{1}{2}x+c \) 和 \( f_2(x)=a_2x^2+bx+\frac{5}{4} \). 且知 \( f_1(-1)=f_2(-1)=f_1(3)=f_2(3)=2 \).

\begin{enumerate}
\item 分别求出产量指标函数 \( f_1(x) \) 和消耗函数 \( f_2(x) \) 的具体表达式；
\item 问因素 \( x \) 取何值时，\( f_1(x) \) 和 \( f_2(x) \) 有最大值或最小值？最大值或最小值各是多少？
\item 画出所求出函数的略图.
\end{enumerate}
\end{problem}

\begin{solution}
\begin{enumerate}
\item 对指标函数 \( f_1(x) \) 来说，由 \( f_1(-1)=f_1(3)=2 \)，可得
\[
\begin{cases}
a_1-\frac{1}{2}+c=2,\\
9a_1+\frac{3}{2}+c=2
\end{cases}
\]
两式相减得 \(8a_1+2=0\)，所以 \(a_1=-\frac14\)；代入第一式，得 \(c=\frac{11}{4}\). 所以
\(f_1(x)=-\frac14x^2+\frac12x+\frac{11}{4}\).

对消耗函数 \( f_2(x) \) 来说，由 \( f_2(-1)=f_2(3)=2 \)，可得
\[
\begin{cases}
a_2-b+\frac{5}{4}=2,\\
9a_2+3b+\frac{5}{4}=2
\end{cases}
\]
两式相减并整理得 \(8a_2+4b=0\)，即 \(b=-2a_2\). 代入第一式得 \(3a_2=\frac34\)，所以 \(a_2=\frac14\)，\(b=-\frac12\). 因此
\(f_2(x)=\frac14x^2-\frac12x+\frac54\).

将所得函数代回 \(x=-1\) 和 \(x=3\)，可得
\(f_1(-1)=f_1(3)=f_2(-1)=f_2(3)=2\)，符合题设条件.

\item 在 \(x\in\R\) 上，\(f_1(x)=-\frac14(x-1)^2+3\)，而 \((x-1)^2\geq0\)，所以 \(f_1(x)\leq3\)，当且仅当 \(x=1\) 时取等号. 因此 \(f_1\) 在 \(x=1\) 时有最大值 \(3\)，且没有最小值.

对 \(f_2\) 配方，得
\[
f_2(x)=\frac14x^2-\frac12x+\frac54
=\frac14\left(x^2-2x+1\right)+1
=\frac14(x-1)^2+1
\]
因为 \((x-1)^2\geq0\)，所以 \(f_2(x)\geq1\)，当且仅当 \(x=1\) 时取等号. 因此 \(f_2\) 在 \(x=1\) 时有最小值 \(1\)，且没有最大值.

\item 两个函数的对称轴都是 \(x=1\). 函数 \(f_1\) 的图象开口向下，顶点为 \((1,3)\)，函数 \(f_2\) 的图象开口向上，顶点为 \((1,1)\). 两个函数的交点满足
\[
-\frac14(x-1)^2+3=\frac14(x-1)^2+1
\]
所以
\(\frac12(x-1)^2=2\)，\((x-1)^2=4\)，\(x=-1\text{ 或 }x=3\).
代入任一函数可得交点的纵坐标为 \(2\)，故交点为 \((-1,2)\) 和 \((3,2)\). 又有
\(f_1(0)=\frac{11}{4}\)，\(f_2(0)=\frac54\)，据此描点并分别作开口向下、向上的抛物线，略图如下：

\end{enumerate}
\solutionfigure{\bitmapfigure[max width=8.0cm]{img/1977/henan/q8_solution_fig1.png}}
\end{solution}

\begin{problem}
已知过点 \( P\left( 0,3\sqrt{2} \right) \) 且斜率为 \( k \) 的直线与圆心在原点，半径为 \( 3 \) 的圆相交于 \(M\)，\(N\) 两点.

\begin{enumerate}
\item 求 \(M\)，\(N\) 的坐标；
\setcounter{enumi}{2}
\item 问当 \(M\)，\(N\) 两点重合时，\( k \) 为何值？过 \(P\) 点的直线与圆的位置关系如何？这样的直线有几条？它们的夹角是多大？
\end{enumerate}
\end{problem}

\begin{solution}
\begin{enumerate}
\item 过点 \(P\left(0,3\sqrt{2}\right)\) 且斜率为 \(k\) 的直线方程是 \(y-3\sqrt{2}=kx\). 圆心在原点、半径为 \(3\) 的圆的方程是 \(x^2+y^2=9\). 交点坐标 \((x,y)\) 应满足方程组
\[
\begin{cases}
y-3\sqrt{2}=kx,\\
x^2+y^2=9
\end{cases}
\]
由 \(y=kx+3\sqrt{2}\)，代入圆的方程，得
\[
(1+k^2)x^2+6\sqrt{2}kx+9=0
\]
由于题设中的 \(M\)，\(N\) 是两个不同的交点，故该关于 \(x\) 的二次方程有两个不等实根，其判别式为
\[
\Delta=(6\sqrt{2}k)^2-4(1+k^2)\cdot9=36(k^2-1)>0
\]
所以 \(k^2>1\). 由求根公式，
\[
x=\frac{-3\sqrt{2}k\pm3\sqrt{k^2-1}}{1+k^2}
\]
再由 \(y=kx+3\sqrt{2}\)，得到
\[
y=\frac{3\sqrt{2}\pm3k\sqrt{k^2-1}}{1+k^2}
\]
故交点 \(M\)，\(N\) 的坐标分别为
\[
\left( \frac{-3\sqrt{2}k+3\sqrt{k^2-1}}{1+k^2},\frac{3\sqrt{2}+3k\sqrt{k^2-1}}{1+k^2} \right)
\]
\[
\left( \frac{-3\sqrt{2}k-3\sqrt{k^2-1}}{1+k^2},\frac{3\sqrt{2}-3k\sqrt{k^2-1}}{1+k^2} \right)
\]

\setcounter{enumi}{2}
\item 当 \(M\)，\(N\) 两点重合时，二次方程的判别式为 \(0\)，故
\(36(k^2-1)=0\)，即 \(k=\pm1\). 此时交点重合，过 \(P\) 点的直线与圆相切. 又因为
\(OP=3\sqrt{2}>3\)，点 \(P\) 在圆外，故这样的切线恰有两条，分别为
\(y=x+3\sqrt{2}\)，\(y=-x+3\sqrt{2}\).
两条切线的斜率分别为 \(1\) 和 \(-1\)，互为负倒数，所以两条切线互相垂直，夹角为 \(90^\circ\).
\end{enumerate}
\end{solution}

\begin{problem}
某电管所为实现农业现代化，加强电力网的建设，沿着一条通往农村的新公路栽电杆. 已知一辆汽车每次从电管所运 \(3\) 根电线杆，相邻两根电线杆相距 \(50\mathrm{~m}\)，汽车往返的总行程是 \(35.5\mathrm{~km}\)，最后一根电线杆与电管所相距 \(2450\mathrm{~m}\).

\begin{enumerate}
\item 第一根电线杆与电管所相距多少？
\item 共栽了多少根电线杆？
\end{enumerate}
\end{problem}

\begin{solution}
\begin{enumerate}
\item 设第一根电线杆与电管所相距 \(x\mathrm{~m}\)，共栽 \(n\) 根. 因为汽车每次运 \(3\) 根，所以往返次数为 \(\frac n3\)，其中 \(n\) 必须是 \(3\) 的倍数. 第一趟把三根电线杆运到距离分别为 \(x\)、\(x+50\)、\(x+100\) 的位置，故往返路程为 \(2(x+100)\mathrm{~m}\). 以后每趟最远的一根比前一趟多 \(150\mathrm{~m}\)，往返路程依次增加 \(300\mathrm{~m}\)，最后一趟路程为 \(2\times2450\mathrm{~m}\). 总行程 \(35.5\mathrm{~km}=35500\mathrm{~m}\).

最后一根电线杆的位置给出
\[
\begin{cases}
2450=x+\left( n-1 \right)50,\\
35500=\frac n3\cdot\frac{2(x+100)+2\times2450}{2}.
\end{cases}
\]
由第一式得 \(x=2500-50n\). 代入第二式，并注意
\(\frac{2(x+100)+2\times2450}{2}=x+2550=50(101-n)\)，可得
\(35500=\frac n3\cdot50(101-n)\)，\(n(101-n)=2130\).
于是
\(n^2-101n+2130=0\)，\((n-30)(n-71)=0\)
所以 \(n=30\) 或 \(n=71\). 当 \(n=30\) 时，\(x=1000\)；当 \(n=71\) 时，\(x=-1050\)，不符合距离为正的实际意义，且 \(71\) 也不是 \(3\) 的倍数，故舍去.

因此第一根电线杆与电管所相距 \(1000\mathrm{~m}\).

\item 共栽了 \(30\) 根电线杆.
\end{enumerate}
\end{solution}

\section{参考题}

\begin{problem}
已知 \( x+x^{-1}=2\cos\theta \)，求证 \( x^n+x^{-n}=2\cos n\theta \).
\end{problem}

\begin{solution}
设 \(n\) 为非负整数，且 \(x\ne0\). 由 \(x+x^{-1}=2\cos\theta\)，两边乘以 \(x\)，得
\(x^2-2\cos\theta x+1=0\). 将其配方可写为
\(\left(x-\cos\theta\right)^2+\sin^2\theta=0\)，所以
\(x=\cos\theta+\i\sin\theta\) 或 \(x=\cos\theta-\i\sin\theta\).

若 \(x=\cos\theta+\i\sin\theta\)，则 \(x^{-1}=\cos\theta-\i\sin\theta\). 由棣莫弗定理，
\[
x^n+x^{-n}
=\left(\cos\theta+\i\sin\theta\right)^n
  +\left(\cos\theta-\i\sin\theta\right)^n
=\left(\cos n\theta+\i\sin n\theta\right)
  +\left(\cos n\theta-\i\sin n\theta\right)
=2\cos n\theta
\]
若 \(x=\cos\theta-\i\sin\theta\)，两项互换，和仍为 \(2\cos n\theta\). 故
\(x^n+x^{-n}=2\cos n\theta\).
\end{solution}

\begin{problem}
\begin{enumerate}
\item 设 \( y=\ln\left( x+\sqrt{x^2-a^2} \right) \)，求 \( y \) 的导数 \( y' \).

\item 求 \( \int_0^{2p}\sqrt{2px}\,\mathrm{d}x \) 的值.
\end{enumerate}
\end{problem}

\begin{solution}
\begin{enumerate}
\item 设 \(a\in\R\). 当 \(a\ne0\) 时，根式要求 \(x\geq|a|\)，且对 \(x\leq-|a|\) 有
\(x+\sqrt{x^2-a^2}<0\)，不在对数函数定义域内，所以函数的定义域为 \([|a|,+\infty)\). 当 \(a=0\) 时，根式为 \(|x|\)，对数真数为 \(x+|x|\)，故定义域为 \((0,+\infty)\). 因而在定义域的内部 \(x>|a|\) 上，由复合函数求导法则，
\[
y'=\frac{1+\dfrac{x}{\sqrt{x^2-a^2}}}{x+\sqrt{x^2-a^2}}
=\frac{\sqrt{x^2-a^2}+x}{\sqrt{x^2-a^2}\left(x+\sqrt{x^2-a^2}\right)}
=\frac{1}{\sqrt{x^2-a^2}}
\]

\item 先考虑 \(p>0\). 令 \(u=2px\)，则 \(\mathrm{d}u=2p\,\mathrm{d}x\)，且积分上下限分别变为 \(0\) 和 \(4p^2\). 因此
\[
\int_0^{2p}\sqrt{2px}\,\mathrm{d}x
=\frac{1}{2p}\int_0^{4p^2}u^{\frac12}\,\mathrm{d}u
=\frac{1}{3p}\left.u^{\frac32}\right|_0^{4p^2}
=\frac{1}{3p}\left(4p^2\right)^{\frac32}
=\frac{8}{3}p^2
\]
若 \(p=0\)，原积分为 \(0\). 若 \(p<0\)，同样令 \(u=2px\)，但此时 \(1/(2p)<0\)，所以
\[
\int_0^{2p}\sqrt{2px}\,\mathrm{d}x
=\frac{1}{2p}\int_0^{4p^2}u^{\frac12}\,\mathrm{d}u
=\frac{1}{3p}\left(4p^2\right)^{\frac32}
=-\frac{8}{3}p^2
\]
因此在 \(p\in\R\) 时，积分值可统一写为 \(\frac83p\lvert p\rvert\)；若题目另有约定 \(p\geq0\)，则结果为 \(\frac83p^2\).
\end{enumerate}
\end{solution}
